\section{Conclusion}
\label{sec:conclusion}

Random-split evaluation makes network intrusion detection look solved: the same classifier that is 99.8\% accurate
on a random split of CIC-IDS2017 loses most of its macro-F1 when it meets the future. We evaluated adaptation
strategies under the conditions of a deployment: chronological replay, attack classes that appear after training,
and a small, delayed labelling budget, with metrics that remain comparable when methods learn different classes.

\nova{} combines class-conditional open-world rejection with per-class recalibration, a novelty-weighted and
diverse query, conservative label propagation, a class-balanced memory and a guarded deployment loop. At a 10\%
budget it is the best or tied-best method on open-world macro-F1 on all three datasets, with a lower false-positive
rate than the strongest baseline on each. At a 1\% budget, it leads on both CIC datasets. Recalibrating the
rejection thresholds of learned classes on recent held-out labels is what keeps it above the never-updated model on
UNSW-NB15, where every other adaptive method falls below it.

The approach has clear limits. Short-lived attacks that resemble benign traffic remain undetected by every method
and every representation we tried, including a meta-learned embedding. The accuracy gain did not carry over to an
MLP classifier, and on one dataset at the smallest budget, open-world scoring favours methods that never learn the
new classes. Future work should evaluate with noisy analysts, longer captures and corrected datasets, and combine
learned and raw-feature detectors to gain the better-calibrated false-positive rate of the former without losing
classes such as infiltration. All runs, configurations and per-flow predictions are released to make such
comparisons cumulative.
